\chapter{Kolmogorov-Komplexität}\label{ch:kolmogorov}
In der algorithmischen Informationstheorie ist die Kolmogorov\footnote{Andrei Nikolajewitsch Kolmogorow, herausragender sowjetischer Mathematiker des 20. Jahrhunderts und Begründer der algorithmischen Komplexitätstheorie}-Komplexität eines Objekts, wie zum Beispiel eines gegebenen Textes, die Länge eines kürzesten Computerprogramms (in einer festgelegten Programmiersprache), das dieses Objekt als Ausgabe (Output) produziert.
\par
Die Kolmogorov-Komplexität gibt uns eine Möglichkeit, den Informationsinhalt von Objekten zu bestimmen, und erlaubt uns, über kürzeste Darstellungen zu sprechen. Zudem ermöglicht sie eine sinnvolle Definition der Zufälligkeit von Texten und Zahlen. \textit{Zufall} ist ein zentraler Begriff der modernen Wissenschaft.

\section{Intuition}
Betrachte die wunderschöne Illustration eines Teils der Mandelbrotmenge in \cref{fig:mandelbrot_16k}.
\begin{figure}[H]
	\centering
	\includegraphics[width=0.6\textwidth]{EF/Kolmogorov/mandelbrot_16k.png}
	\caption{Mandelbrotmenge in 16K-Auflösung}
	\label{fig:mandelbrot_16k}
\end{figure}
Eine hochauflösende Rastergrafik muss die Farbinformationen sehr vieler Pixel speichern und kann daher selbst als komprimierte PNG-Datei rund $100$ Millionen Bits Speicherplatz benötigen. Ein recht kurzes Computerprogramm (siehe \cref{listing:mandelbrot}) kann jedoch eine solche Darstellung erzeugen, indem es die Definition der Mandelbrotmenge verwendet. Das Programm ist damit eine deutlich kürzere Beschreibung des Bildmotivs als die Angabe jedes Pixels. Für die Erzeugung der Darstellung benötigt es zahlreiche Rechenoperationen und viel Rechenzeit. Dies ändert aber nichts daran, dass es sich um eine kurze Beschreibung der Mandelbrotgrafik handelt. An dieser Stelle können wir wärmstens empfehlen, einen Blick in das Werk\cite{LowComplexityArt} zu werfen.

\begin{myexercise}
		Betrachte die folgenden zwei Bilder. \cref{fig:fig1} zeigt ein Muster aus zahlreichen Kreisen, während \cref{fig:fig2} zufälliges Rauschen (White Noise) zeigt.

	\begin{figure}[H]
		\centering
		\begin{minipage}[b]{0.45\textwidth}
			\centering
			\includegraphics[width=\textwidth]{EF/Kolmogorov/circles.pdf}
			\caption{Kreise}
			\label{fig:fig1}
		\end{minipage}
		\hspace{0.05\textwidth}
		\begin{minipage}[b]{0.45\textwidth}
			\centering
			\includegraphics[width=\textwidth]{EF/Kolmogorov/white_noise.pdf}
			\caption{Rauschen}
			\label{fig:fig2}
		\end{minipage}
	\end{figure}

	Welches dieser beiden Bilder erlaubt vermutlich eine kürzere Beschreibung? Begründe deine Antwort.
\begin{myanswer}
	Das Kreismuster in \cref{fig:fig1} erlaubt vermutlich eine deutlich kürzere Beschreibung. Es ist sehr regelmässig und kann durch ein kurzes Computerprogramm beschrieben werden, das in einer Schleife Kreise an bestimmten Positionen zeichnet (wir erinnern uns an die Python-Turtle). Ein tatsächlich gleichverteilt zufällig erzeugtes Rauschbild ist hingegen mit hoher Wahrscheinlichkeit kaum komprimierbar: Ein erzeugendes Programm muss dann im Wesentlichen die Farbinformation (weiss beziehungsweise schwarz) jedes einzelnen Pixels speichern. Vom sichtbaren Eindruck allein lässt sich dies allerdings nicht beweisen, denn ein ähnlich aussehendes Bild könnte auch von einem kurzen Pseudozufallszahlengenerator mit einem kurzen Startwert erzeugt worden sein.
\end{myanswer}
\end{myexercise}

\begin{myremark}
	In diesem Kapitel fassen wir Wörter (Texte) als Träger von Information auf. Obwohl das Konzept der Kolmogorov-Komplexität auf beliebige Objekte anwendbar ist, werden wir uns hier der Einfachheit halber nur mit binären Wörtern befassen.
\end{myremark}

Wir werden nun allmählich versuchen, eine robuste Methode zur Messung des Informationsgehalts eines Wortes zu entwickeln und dabei mehrfach auf \cref{ch:Details} verweisen.

\begin{myexample}\label{myexample:passwort}
Alice hat ein langes binäres Passwort für den Zugriff auf eine Streamingplattform gewählt.
Bob möchte gerne die neueste Staffel seiner Lieblingsserie schauen und ruft dazu Alice an, um sie um das Passwort zu bitten\footnote{Bob scheint nicht mit dem \textit{BitTorrent}-Kommunikationsprotokoll (Peer-to-Peer-Filesharing) vertraut zu sein.
	Vielleicht benötigt er aber auch nur einen Vorwand, um Alice anzurufen.}.

\begin{itemize}
	\item Angenommen, Alices Passwort lautet $p_0 := 01011010110001101101111010$. Wie könnte sie Bob das Passwort $p_0$ per Telefon mitteilen? Ihr wird (mehr oder weniger) nichts anderes übrig bleiben, als Bob das Passwort Bit um Bit (Ziffer um Ziffer) zu diktieren. $p_0$ lässt sich also nicht kurz beschreiben und man kann sich $p_0$ auch nicht einfach merken.
	\item Anders hingegen wäre die Situation bei dem Passwort
	      \begin{align*}
		      p_1 := 11100111001110011100111001110011100 = (11100)^7.
	      \end{align*}
	      Das Passwort $p_1$ weist eine starke Regelmässigkeit auf. Natürlich könnte Alice auch dieses Passwort Bit für Bit an Bob diktieren. Offensichtlich ist es jedoch einfacher, wenn sie ihm einfach mitteilt:
	      \begin{center}
		      \enquote{Du erhältst das Passwort, indem du siebenmal das Wort $11100$ hintereinander schreibst.}
	      \end{center}
	      Dadurch hat Alice die Darstellung von $p_1$ \textit{komprimiert}. Bitte beachte, dass Bob sich dieses Passwort $p_1$ recht einfach merken könnte.
\end{itemize}
\end{myexample}

\cref{myexample:passwort} motiviert die folgende intuitive Vorstellung: Wir wollen einem Wort einen kleinen Informationsgehalt beimessen, falls es eine kurze Darstellung/Beschreibung besitzt (komprimierbar ist), und
einen grossen Informationsgehalt, falls das Wort so unregelmässig ist, dass es keine kurze Darstellung/Beschreibung zulässt.

\begin{mydefinition}[Komprimierung intuitiv]\label{mydefinition:Komprimierung}
Eine kürzere Darstellung eines Wortes $w$, aus der $w$ durch ein festgelegtes Verfahren eindeutig rekonstruiert werden kann, nennen wir eine verlustfreie Komprimierung von $w$.
\end{mydefinition}

\begin{myexercise}\label{myexercise:wahrscheinlichkeit}
	Es sei $S$ die Menge aller binären Wörter der Länge $9$. Nun \enquote{ziehst} du gemäss der Gleichverteilung ein Wort aus $S$.
	\begin{enumerate}
		\item Mit welcher Wahrscheinlichkeit wirst du das Wort $010010110$ ziehen?
		\item Mit welcher Wahrscheinlichkeit wirst du das Wort $111111111$ ziehen?
	\end{enumerate}
	\begin{myanswer}
		Die Menge $S$ enthält genau $2^9 = 512$ verschiedene binäre Wörter der Länge $9$.
		\begin{enumerate}
			\item Die Wahrscheinlichkeit, das Wort $010010110$ zu ziehen, ist $\frac{1}{512}$.
			\item Die Wahrscheinlichkeit, das Wort $111111111$ zu ziehen, ist ebenfalls $\frac{1}{512}$.
		\end{enumerate}
	\end{myanswer}
\end{myexercise}

\begin{myexercise}
	Intuitiv sollte \textit{zufällig} bedeuten: \enquote{nach keinem klaren Plan gebaut}. Ein zufälliges Objekt weist also eine chaotische (und keine regelmässige) Struktur auf. Vergleiche mit \cref{myexercise:wahrscheinlichkeit} und argumentiere, warum die klassische Wahrscheinlichkeitsrechnung diese Intuition nicht einfängt.
\begin{myanswer}
	Wie in \cref{myexercise:wahrscheinlichkeit} gesehen, ist aus Sicht der klassischen Wahrscheinlichkeitstheorie jedes binäre Wort der Länge $9$ gleich wahrscheinlich. Die klassische Wahrscheinlichkeitstheorie macht also keine Aussage über die Zufälligkeit eines einzelnen Wortes, sondern nur über die Wahrscheinlichkeit, mit der ein Wort aus einer Menge von Wörtern gezogen wird.
	
	Obwohl das Wort $111111111$ sehr regelmässig und das Wort $010010110$ sehr unregelmässig erscheint, sind beide unter der Gleichverteilung gleich wahrscheinlich. Die Wahrscheinlichkeit eines einzelnen Ergebnisses unterscheidet hier also nicht zwischen strukturierten und unstrukturierten Wörtern. Die Kolmogorov-Komplexität ermöglicht eine solche Unterscheidung anhand der kürzesten algorithmischen Beschreibung.
\end{myanswer}
\end{myexercise}

Ein denkbares Vorgehen zur Messung des Informationsgehalts eines binären Wortes $w$ wäre, sich auf eine konkrete Komprimierungsmethode $komp$ zu einigen.
Die Länge $\abs{komp(w)}$ des komprimierten Wortes $komp(w)$ könnte dann als Mass für den Informationsgehalt von $w$ angesehen werden.
\par
Natürlich muss auch $komp(w)$ ein Wort über dem binären Alphabet sein.
Würden wir $komp(w)$ in einem anderen Alphabet darstellen, wäre der Vergleich der Darstellungslängen nicht mehr sinnvoll, denn schliesslich ist es immer möglich, ein Wort in einem mächtigen Alphabet kurz darzustellen.

\section{Problematik einer fest gewählten Komprimierung}
Es gibt unendlich viele denkbare Komprimierungsmethoden.
Doch welche Komprimierungsmethode stellt nun die richtige Wahl dar? Unabhängig von der gewählten Methode gibt es stets eine andere Komprimierung, die für unendlich viele Wörter kürzere Darstellungen erzeugt.

\begin{myexample}[Komprimierung durch Ausnutzung von Wiederholungen]\label{myexample:Wiederholungen}
Lass uns ein binäres Wort mit zahlreichen Wiederholungen von Teilwörtern untersuchen. Konkret wollen wir das Wort
\begin{align*}
	w := 00(101)^{\textcolor{Red}{25}}(01010)^{\textcolor{Blue}{9}}(1011)^{\textcolor{Green}{16}} = 00(101)\textcolor{Red}{11001}(01010)\textcolor{Blue}{1001}(1011)\textcolor{Green}{10000}
\end{align*}
untersuchen, wobei wir die Exponenten in Basis 10 durch ihre binäre Darstellung ersetzt haben. Die resultierende Darstellung ist jedoch noch immer kein binäres Wort. Tatsächlich ist es ein Wort über dem Alphabet $\lrc{0, 1, (, )}$, das vier Symbole enthält. Mit der Vereinbarung
\begin{align*}
	0 \to 00, \quad 1 \to 11, \quad ( \to 10, \quad ) \to 01
\end{align*}
erhalten wir die doppelt so lange (aber dafür binäre) Darstellung:
\begin{align*}
	w = 00001011001101111100001110001100110001110000111011001111011100000000.
\end{align*}
\end{myexample}

\begin{myexample}[Komprimierung durch Primfaktorisierung]\label{myexample:Primkomprimierung}
Es sei $N\geq2$ eine natürliche Zahl und $x:=bin(N)$ ihre Binärdarstellung ohne führende Nullen. Dann gilt
\begin{align*}
	N=val_2(x) = p_1^{e_1}p_2^{e_2}\ldots p_t^{e_t},
\end{align*}
wobei die $p_i$ unterschiedliche Primzahlen sind und $e_i\geq 1$ für $i\in\setcm{n\in\N}{1\leq n\leq t}$.
Eine denkbare Darstellung von $p_1^{e_1}p_2^{e_2}\ldots p_t^{e_t}$ ist
\begin{align*}
	bin(p_1)(bin(e_1))bin(p_2)(bin(e_2))\ldots bin(p_t)(bin(e_t)).
\end{align*}
Die Klammern trennen die einzelnen Bestandteile eindeutig. Mit der Vereinbarung aus \cref{myexample:Wiederholungen} erhalten wir wieder eine eindeutig dekodierbare binäre Darstellung von $x$.
\end{myexample}

Wir bezeichnen mit $L_W(x)$ die Länge der Darstellung von $x$, die durch das Ausnutzen von Wiederholungen entsteht, und mit $L_P(x)$ die Länge der Darstellung, die mithilfe der Primfaktorisierung entsteht. Die beiden Methoden nutzen unterschiedliche Arten von Regelmässigkeit aus: Die erste Methode ist besonders erfolgreich, wenn ein Wort aus langen Wiederholungen kurzer Teilwörter besteht. Die zweite Methode ist besonders erfolgreich, wenn der Zahlenwert eines Wortes das Produkt hoher Potenzen weniger Primzahlen ist. Daher ist nicht zu erwarten, dass eine der beiden Methoden für alle Wörter mindestens so kurze Darstellungen wie die andere erzeugt. In \cref{myexample:Unvergleichbarkeit} illustrieren wir diese unterschiedlichen Stärken anhand zweier Beispiele.

\begin{myexample}[Unterschiedliche Stärken von Komprimierungsmethoden]\label{myexample:Unvergleichbarkeit}
	\leavevmode
\begin{itemize}
	\item Für das Wort $w_A := (10110010)^{512}$ mit der binären Länge von $8 \cdot 512 = 4096$ Bits, ergibt die Komprimierung aus \cref{myexample:Wiederholungen}:
\begin{align*}
	w_A &= (10110010)^{512} = (10110010)bin(512) = (10110010)1000000000 =\\
	&= \underbrace{10}_{(}\underbrace{1100111100001100}_{10110010}\underbrace{01}_{)}\underbrace{11000000000000000000}_{1000000000}
\end{align*}
also als zusammenhängendes binäres Wort:
\begin{align*}
	1010110011110000110001011100000000000000
\end{align*}
mit einer Darstellungslänge von $40$ Bits gegenüber der unkomprimierten Darstellungslänge von $4096$ Bits.
	\item Für das Wort $w_B := bin\lr{3^{5000}\cdot 5^{4000}\cdot 17^{3000}}$ mit der binären Länge von $17230$ Bits (siehe \cref{mytheorem:AnzahlZiffern}), ergibt die Komprimierung aus \cref{myexample:Primkomprimierung}:
\begin{align*}
	w_B &= bin\lr{3^{5000}\cdot 5^{4000}\cdot 17^{3000}} =\\
	&= bin(3)(bin(5000))bin(5)(bin(4000))bin(17)(bin(3000)) =\\
	&= 11(1001110001000)101(111110100000)10001(101110111000).
\end{align*}
Nach Anwendung unserer Vereinbarung
\begin{align*}
	0 \to 00, \quad 1 \to 11, \quad ( \to 10, \quad ) \to 01
\end{align*}
erhalten wir eine Darstellungslänge von $106$ Bits gegenüber der unkomprimierten Darstellungslänge von $17230$ Bits.
	\item Heuristisch erwarten wir, dass die Primfaktorisierung für $w_A$ kaum gewinnbringend ist, da $val_2(w_A)$ vermutlich nicht als Produkt hoher Potenzen weniger Primzahlen geschrieben werden kann. Umgekehrt erwarten wir bei $w_B$ keine Wiederholungsmuster, deren Ausnutzung sich lohnt. Diese Überlegungen illustrieren die unterschiedlichen Stärken der Methoden, beweisen aber keine formale Unvergleichbarkeit. Da sowohl $w_A$ als auch $w_B$ mit einer $1$ beginnen, geht bei der Interpretation als Zahlenwert keine Information über führende Nullen verloren.
\end{itemize}
\end{myexample}

Die Beispiele illustrieren, dass verschiedene Komprimierungsmethoden unterschiedliche Strukturen ausnutzen. Keine der beiden Methoden erfasst offensichtlich alle Arten von Regelmässigkeit.

\section{Idee hinter der Kolmogorov-Komplexität}
Die Definition eines Komplexitätsmasses sollte in dem Sinne robust sein, dass sie nicht auf der willkürlichen Wahl einer Komprimierungsmethode basiert. Ein Komplexitätsmass ohne eine solche willkürliche Wahl ist die Kolmogorov-Komplexität. Um deren Definition zu verstehen, müssen wir zunächst den Begriff der \textit{Generierung} eines Wortes einführen.

\begin{mydefinition}[Generierung eines Wortes]\label{mydefinition:erzeugung}
Es sei $w$ ein beliebiges Wort. Wir sagen, dass ein Programm (ein Algorithmus) $A$ das Wort $w$ generiert, falls der Aufruf $A()$ genau das Wort $w$ ausgibt.
\end{mydefinition}

\begin{myexample}\label{myexample:generateA}
Das Programm
\begin{lstlisting}[language=C++]
#include <iostream>
void A() {
    std::cout << "00110111";
}
\end{lstlisting}
generiert das Wort $00110111$.
\end{myexample}

\begin{myexample}
Das Programm
\begin{lstlisting}[language=C++]
#include <iostream>
#include <string>
void B() {
    std::string wort(9900507, '1');
    std::cout << wort;
}
\end{lstlisting}
generiert das Wort $w := 1^{9900507} = \underbrace{11\ldots 111}_{\text{9900507 Einsen}}$.
\par
Beachte, dass der Algorithmus \texttt{B} eine kurze Beschreibung von $w$ relativ zur grossen Länge von $w$ darstellt.
\end{myexample}

Die Idee von Kolmogorov besteht darin, dass wir uns nicht auf eine willkürlich gewählte Komprimierung einigen müssen. Stattdessen erlaubt er beliebige Programme (in einer fest gewählten Programmiersprache) zur Beschreibung von Texten. Wenn ein Text durch ein kurzes Programm beschrieben werden kann, so besitzt er eine geringe Kolmogorov-Komplexität. Erlaubt ein Text hingegen keine kurze Beschreibung als Programm, so besitzt er eine hohe Kolmogorov-Komplexität.

\section{Binäre Datenblöcke und Definition der Kolmogorov-Komplexität}
Zunächst möchten wir klären, wie ein binäres Wort in einem Programm gespeichert werden kann, ohne dass jedes seiner Bits als gewöhnliches Textzeichen kodiert werden muss. Betrachten wir dazu nochmals das folgende Standard-\texttt{C++}-Programm aus \cref{myexample:generateA}:
\begin{lstlisting}[language=C++,escapechar=!] 
#include <iostream>
void A() {
    std::cout << "00110111";
}
\end{lstlisting}

Das Problem ist, dass wir typischerweise die Symbole der Tastatur durch Folgen von Nullen und Einsen darstellen müssen (zum Beispiel mittels der ASCII-Kodierung). Wenn wir das gesamte Programm \texttt{A} auf diese Weise kodieren, benötigt jedes als Textzeichen notierte Bit von $w := 00110111$ mehrere Bits im Maschinencode. Diesen konstanten Faktor möchten wir vermeiden. Dazu entwerfen wir eine leicht modifizierte Version des \texttt{g++}-Compilers und der Programmiersprache \texttt{C++}. Wir nennen den modifizierten Compiler \texttt{g'++} und die angepasste Sprache \texttt{C'++}.
\par
Die Sprache \texttt{C'++} unterscheidet sich von Standard-\texttt{C++} darin, dass in einem \texttt{C'++}-Programm binäre Blöcke durch die Escape-Sequenzen \cppinline{***} und \cppinline{+++} markiert werden dürfen. Zwischen diesen Begrenzungszeichen steht ein binäres Wort $w$. Kommen die Escape-Sequenzen in einem Programm mehrfach vor, so muss zwischen ihnen stets dasselbe Wort $w$ stehen; andernfalls meldet \texttt{g'++} einen semantischen Fehler.
\par
Der Compiler \texttt{g'++} speichert diesen gemeinsamen Block $w$ genau einmal und unverändert als binären Rohdatenblock im erzeugten Maschinencode. Alle Vorkommen von \cppinline{***w+++} werden durch Verweise auf diesen einen Block ersetzt. Der Rohdatenblock wird am Ende des Maschinencodes abgelegt; der gewöhnliche Programmteil ist eindeutig kodiert, sodass der gesamte verbleibende Teil des Maschinencodes als $w$ erkannt wird. Die Begrenzungszeichen selbst werden nicht in den Maschinencode übernommen. Somit trägt $w$ unabhängig von der Anzahl seiner Vorkommen im Quelltext genau $\abs{w}$ Bits zur Länge des Maschinencodes bei. Die Verweise und der übrige Programmcode werden bei der Länge des Maschinencodes selbstverständlich mitgezählt.
\par
In den folgenden Programmschemata stehen Ausdrücke wie \cppinline{***n+++} beziehungsweise \cppinline{***n**2+++} dafür, dass für jedes feste $n$ konkret die Binärdarstellung von $n$ beziehungsweise von $n^2$ zwischen den Begrenzungszeichen steht. Die Ausdrücke \cppinline{n} und \cppinline{n**2} selbst sind also keine Bestandteile des kompilierten Programms.

\begin{myexample}
	Das Programm
		\begin{lstlisting}[language=C++,escapechar=!] 
        !\colorbox{Yellow}{\#include <iostream>}!
        !\colorbox{Yellow}{void A() \{}!
            !\colorbox{Yellow}{\hspace*{1.5em}std::cout << \texttt{"}***}!!\colorbox{YellowGreen}{00110111}!!\colorbox{Yellow}{+++\texttt{"};}!
        !\colorbox{Yellow}{\texttt{\}}}!
        \end{lstlisting}
generiert das Wort $w := 00110111$ und ist ein gültiges Programm in der Sprache \texttt{C'++}. Der Programmteil \colorbox{YellowGreen}{00110111} wird bei der Kompilierung mithilfe von \texttt{g'++} einmal unverändert im Maschinencode gespeichert und trägt somit exakt $8$ Bits zur Länge des Maschinencodes bei.
\end{myexample}

Nun können wir endlich die formale Definition der Kolmogorov-Komplexität angeben:
\begin{mydefinition}[Kolmogorov-Komplexität]\label{mydefinition:kolmogorov}
Es sei $w$ ein binäres Wort. Die plain Kolmogorov-Komplexität $C(w)$ von $w$ ist definiert als das Minimum der \textbf{binären Längen} aller mit \texttt{g'++} kompilierten \texttt{C'++}-Programme, die $w$ generieren. Im Folgenden nennen wir diese Grösse kurz Kolmogorov-Komplexität.
\end{mydefinition}
Für ein binäres Wort $w$ betrachten wir alle Maschinencodes der \texttt{C'++}-Programme, die $w$ generieren. Die Länge eines kürzesten\footnote{Im Allgemeinen kann es mehrere verschiedene kürzeste Maschinencodes geben, die $w$ generieren.} solchen Maschinencodes ist dann die Zahl $C(w)$. Die Notation $C$ macht deutlich, dass wir die plain Kolmogorov-Komplexität und nicht die präfixfreie Kolmogorov-Komplexität betrachten.

\begin{myremark}
	In \cref{sec:Sprachinvarianz} werden wir sehen, dass die Wahl der Programmiersprache keine wesentliche Rolle spielt.
\end{myremark}
\begin{myremark}
	Ist $C(w)$ ein geeignetes Mass für den Informationsgehalt des binären Wortes $w$? Ja, denn die Kolmogorov-Komplexität berücksichtigt jede berechenbare verlustfreie Komprimierungsmethode. Es sei $x$ ein binäres Wort und $f$ eine solche Komprimierungsmethode, die zu $x$ eine komprimierte Darstellung $x':=f(x)$ erzeugt. Dann existiert ein Dekompressionsalgorithmus $d$ mit $d(f(x))=d(x')=x$. Wir können ein Programm schreiben, das lediglich das kurze Wort $x'$ als Parameter enthält, $x$ mithilfe von $d$ rekonstruiert und anschliessend ausgibt. Der feste Algorithmus $d$ muss zwar ebenfalls im Programm kodiert sein, seine Beschreibungslänge ist jedoch unabhängig von $x$ und $x'$ und kann daher als Konstante betrachtet werden.
\end{myremark}

\section{Invarianz der Programmiersprache}\label{sec:Sprachinvarianz}
Man könnte denken, dass die Wahl einer festen Programmiersprache wiederum ein willkürliches Element einbringt. Dem ist aber nicht so! Wir klären diesen Sachverhalt in \cref{mytheorem:Sprachinvarianz}.

\begin{mydefinition}[Optimale universelle Beschreibungssprache]
	Eine berechenbare Beschreibungssprache $U$ heisst \textit{universell}, wenn sie jedes Programm jeder berechenbaren Beschreibungssprache $S$ mithilfe eines festen Interpreters für $S$ ausführen kann. Sie heisst \textit{optimal}, wenn das Hinzufügen dieses Interpreters die Beschreibungslänge nur um eine von $S$, nicht aber vom beschriebenen Wort abhängige Konstante vergrössert. Wir nehmen an, dass die durch \texttt{C'++} und \texttt{g'++} festgelegte Beschreibungssprache optimal universell ist.
\end{mydefinition}

\begin{mytheorem}[Invarianzsatz]\label{mytheorem:Sprachinvarianz}
	Es sei $U$ eine optimale universelle Beschreibungssprache und $S$ eine beliebige berechenbare Beschreibungssprache. Mit $C_U(w)$ beziehungsweise $C_S(w)$ bezeichnen wir die Länge einer kürzesten Beschreibung von $w$ in der jeweiligen Sprache. Dann existiert eine nur von $S$ abhängige Konstante $c_S$, sodass
	\begin{align*}
		C_U(w)\leq C_S(w)+c_S
	\end{align*}
	für alle binären Wörter $w$ gilt.
	\begin{myproof}
		Da $U$ universell ist, besitzt sie für die feste Sprache $S$ einen Interpreter mit einer festen binären Kennung $i_S$. Für jedes $S$-Programm $p$ gilt
		\begin{align*}
			U(i_Sp)=S(p).
		\end{align*}
		Ist $p$ eine kürzeste $S$-Beschreibung von $w$, so ist $i_Sp$ eine $U$-Beschreibung von $w$. Daher gilt
		\begin{align*}
			C_U(w)\leq \abs{i_S}+\abs{p}=C_S(w)+c_S
		\end{align*}
		mit der von $w$ unabhängigen Konstante $c_S:=\abs{i_S}$.
	\end{myproof}
\end{mytheorem}
Sind $U$ und $V$ beide optimale universelle Beschreibungssprachen, wenden wir den Satz in beide Richtungen an und erhalten
\begin{align*}
	\abs{C_U(w)-C_V(w)}\leq c_{U,V}
\end{align*}
für eine von $w$ unabhängige Konstante $c_{U,V}$. In diesem präzisen Sinne spielt die Wahl der optimalen universellen Beschreibungssprache keine wesentliche Rolle.

\section{Triviale obere Schranke für die Kolmogorov-Komplexität}
Ein beliebiges binäres Wort $w$ kann sicherlich immer beschrieben werden, indem man jedes Bit des Wortes einzeln angibt, $w$ also Bit für Bit diktiert. Diese Beobachtung lässt vermuten, dass die Kolmogorov-Komplexität eines jeden binären Wortes $w$ zumindest nicht (wesentlich) grösser ist als die Anzahl Bits von $w$ (Länge von $w$). Diese Vermutung ist korrekt, wie \cref{mytheorem:upperbound} zeigt.

\begin{mytheorem}[$C(w)$ ist auf keinen Fall wesentlich länger als $\abs{w}$]\label{mytheorem:upperbound}
	Es existiert eine Konstante $c$, sodass für jedes binäre Wort $w$ gilt
	\begin{align*}
		C(w) \leq \abs{w} + c.
	\end{align*}
	\begin{myproof}
		Es genügt, ein Programm anzugeben, das ein beliebiges binäres Wort $w$ generiert und dessen binäre Länge höchstens $\abs{w} + c$ beträgt (für eine geeignete Konstante $c$). Das Programm
		\begin{lstlisting}[language=C++,escapechar=!] 
        !\colorbox{Yellow}{\#include <iostream>}!
        !\colorbox{Yellow}{void A() \{}!
            !\colorbox{Yellow}{\hspace*{1.5em}std::cout << \texttt{"}***}!!\colorbox{YellowGreen}{w}!!\colorbox{Yellow}{+++\texttt{"};}!
        !\colorbox{Yellow}{\texttt{\}}}!
        \end{lstlisting}
		generiert $w$. Der \colorbox{Yellow}{gelbe} Teil des Programms ist für jedes Wort $w$ identisch. Das Wort \colorbox{YellowGreen}{\texttt{w}} kommt unverändert (binär) im Maschinencode vor. Der resultierende Maschinencode hat also eine Länge von $\abs{w}$ Bits (für die Darstellung des binären Wortes $w$) und zusätzlich eine konstante Anzahl Bits $c$ für die Kodierung des eigentlichen Programms (gelber Teil). Die Anzahl Bits $c$ ist in dem Sinne konstant, dass sie nicht von $w$ abhängt. Damit haben wir ein Programm angegeben, das $w$ generiert und eine binäre Länge von $\abs{w} + c$ hat. Somit gilt $C(w) \leq \abs{w} + c$ für eine Konstante $c$.
	\end{myproof}
\end{mytheorem}
Die Kolmogorov-Komplexität eines Wortes $w$ ist also sicherlich (bis auf eine Konstante) nicht grösser als die Länge von $w$.

\section{Sehr regelmässige Wörter}
\cref{mytheorem:upperbound} gibt lediglich eine obere Schranke für die Kolmogorov-Komplexität von binären Wörtern an. Wörter $w$, die eine hohe Regelmässigkeit aufweisen, müssen nicht Bit für Bit angegeben werden, sondern lassen sich (unter Ausnutzung ihrer Regelmässigkeit) deutlich kürzer beschreiben als mit $\abs{w} + c$ vielen Bits.

\begin{myexample}\label{myexample:wn}
Es sei $w_n := 1^n$ für eine beliebige natürliche Zahl $n > 0$. Das Wort $w_n$ besteht also genau aus $n$ Einsen (und hat insbesondere die Länge $n$). Es handelt sich somit um ein sehr regelmässiges Wort, das (intuitiv) nicht Bit für Bit beschrieben werden muss. Das folgende Programm \texttt{Cn} generiert das Wort $w_n$:
\begin{lstlisting}[language=C++,escapechar=!] 
#include <iostream>
#include <string>

void Cn() {
    std::string wort(***!\colorbox{YellowGreen}{n}!+++, '1');
    std::cout << wort;
}
\end{lstlisting}
Im Programm \texttt{Cn} ist einzig die binäre Kodierung \colorbox{YellowGreen}{\texttt{n}} der Zahl $n$ abhängig von $n$ beziehungsweise von $w_n$. Der Rest des Programms ist unabhängig von der konkreten Wahl von $n$ (für jede Wahl von $n$ gleich). Gemäss \cref{mytheorem:AnzahlZiffern} benötigen wir zur binären Darstellung von $n$ genau $\ceil{\log_2(n+1)}$ viele Bits. Damit existieren Konstanten $c_0$ und $c_1$, sodass
\begin{align*}
	C(w_n) \leq c_0 + \ceil{\log_2(n+1)} \stackrel{\cref{mytheorem:Estimate}}{\leq} \log_2(n) + c_1 = \log_2(\abs{w_n}) + c_1
\end{align*}
für beliebiges natürliches $n>0$.
\end{myexample}

\begin{myexample}\label{myexample:vn}
Es sei $v_n := 1^{(n^2)}$ für eine beliebige natürliche Zahl $n > 0$. Das Wort $v_n$ besteht also genau aus $n^2$-vielen Einsen. Es gilt $\abs{v_n} = n^2$ und somit $n = \sqrt{\abs{v_n}}$. Das folgende Programm \texttt{Dn} generiert das Wort $v_n$:
\begin{lstlisting}[language=C++,escapechar=!] 
#include <iostream>
#include <string>

void Dn() {
    std::string wort;
    Natural M = ***!\colorbox{YellowGreen}{n}!+++;
    M = M * M;
    wort.reserve(M);

    for (int i = 0; i < M; ++i) {
        wort += '1';
    }

    std::cout << wort;
}
    \end{lstlisting}
Analog zu \cref{myexample:wn} ist auch hier einzig die binäre Kodierung \colorbox{YellowGreen}{\texttt{n}} der Zahl $n$ abhängig von $n$ beziehungsweise von $v_n$. Der Rest des Programms ist wieder unabhängig von der konkreten Wahl von $n$ (für jede Wahl von $n$ gleich). Zur binären Darstellung von $n$ benötigen wir genau $\ceil{\log_2(n+1)}$ viele Bits. Damit existieren Konstanten $c_3$ und $c_4$, sodass
\begin{align*}
	C(v_n) \leq c_3 + \ceil{\log_2(n+1)} \stackrel{\cref{mytheorem:Estimate}}{\leq} \log_2(n) + c_4 = \log_2\lr{\sqrt{\abs{v_n}}} + c_4
\end{align*}
für beliebiges natürliches $n>0$.

Bei der Bestimmung einer oberen Schranke für die Kolmogorov-Komplexität eines Wortes interessieren wir uns lediglich dafür, wie viele Bits wir für die \textbf{Beschreibung} dieses Wortes benötigen!
Dass
\begin{itemize}
	\item die Inhalte von Variablen wie zum Beispiel von \texttt{M} während der Ausführung des Programms \texttt{Dn} möglicherweise gigantisch gross werden,
	\item die Programmausführung sehr viel Speicher belegen könnte,
	\item oder eine Berechnung wie \texttt{M * M} potenziell sehr viel Rechenzeit in Anspruch nehmen könnte
\end{itemize}
tut hier nichts zur Sache! Uns interessiert einzig die binäre Länge der \textbf{Beschreibung} eines Wortes. Die Programme in diesem Abschnitt sind Pseudocode; der Typ \cppinline{Natural} kann beliebig grosse natürliche Zahlen speichern. Ausserdem nehmen wir an, dass genügend Arbeitsspeicher vorhanden ist.
\end{myexample}

\begin{myexercise}
	Betrachte nochmals \cref{myexample:vn}. Natürlich könnte das Wort $v_n$ auch durch das folgende Programm beschrieben werden:
	\begin{lstlisting}[language=C++,escapechar=!] 
#include <iostream>
#include <string>

void Fn() {
    std::string wort;
    wort.reserve(***n**2+++);

    for (int i = 0; i < ***n**2+++; ++i) {
        wort += '1';
    }

    std::cout << wort;
}
    \end{lstlisting}
	Warum beweist dieses Programm sicherlich eine schlechtere obere Schranke für die Kolmogorov-Komplexität von $v_n$ als \cref{myexample:vn}?
	\begin{myanswer}
		Die beiden Vorkommen von \cppinline{***n**2+++} bezeichnen denselben binären Block. Der Compiler \texttt{g'++} speichert die Binärdarstellung von $n^2$ deshalb nur einmal. Damit liefert \texttt{Fn} die obere Schranke
		\begin{align*}
			C(v_n)\leq c + \ceil{\log_2\lr{n^2 + 1}} \leq c' + \log_2\lr{n^2} = c' + \log_2\lr{\abs{v_n}}
		\end{align*}
		für geeignete Konstanten $c, c'$. Diese Schranke ist schlechter als die in \cref{myexample:vn} hergeleitete Schranke
		\begin{align*}
			C(v_n)\leq c_4 + \log_2\lr{\sqrt{\abs{v_n}}},
		\end{align*}
		da \texttt{Fn} die längere Binärdarstellung von $n^2$ speichert, während \texttt{Dn} nur die Binärdarstellung von $n$ speichert und $n^2$ zur Laufzeit berechnet.
	\end{myanswer}
\end{myexercise}

\section{Nichtkomprimierbare Wörter und Zufall}
Wir haben die Kolmogorov-Komplexität nur für binäre Wörter definiert. Wir können die Kolmogorov-Komplexität aber auch ganz einfach für natürliche Zahlen definieren, indem wir die natürliche Zahl einfach in Basis 2 (binär) darstellen, um ein binäres Wort zu erhalten.

\begin{mydefinition}[Kolmogorov-Komplexität einer natürlichen Zahl]
Die Kolmogorov-Komplexität $C(n)$ einer natürlichen Zahl $n$ ist $C(n) := C(bin(n))$.
\end{mydefinition}

Es existieren Wörter, die nicht komprimierbar sind, die also eine Kolmogorov-Komplexität besitzen, die nicht kleiner ist als ihre Länge.

\begin{mytheorem}[Existenz nichtkomprimierbarer Wörter]\label{mytheorem:nichtkomprimierbar}
	Für jede natürliche Zahl $n>0$ existiert mindestens ein binäres Wort $w$ der Länge $n$, für das gilt
	\begin{align*}
		C(w)\geq n = \abs{w},
	\end{align*}
	das heisst, für jede Länge $n$ existiert mindestens ein Wort dieser Länge, das sich nicht komprimieren lässt.

	\vspace{\baselineskip}
	Insbesondere ist dadurch für jedes natürliche $n>0$ die Existenz mindestens eines binären Wortes $w$ (wir sagen nichts über die Länge von $w$ aus) mit $C(w)\geq n$ gesichert.
	\begin{myproof}
		Unser Beweis verwendet ein einfaches kombinatorisches Argument (Abzählen): Ein Maschinencode kann höchstens ein binäres Wort generieren. Die Anzahl aller Maschinencodes der Länge kleiner als $n$ ist
		\begin{align*}
			\sum_{k=0}^{n-1}2^k=2^n-1.
		\end{align*}
		Es gibt jedoch genau $2^n$ binäre Wörter der Länge $n$. Somit existiert mindestens ein binäres Wort $x$ der Länge $n$, das von keinem Maschinencode der Länge $<n$ generiert wird. Daher gilt $C(x)\geq n=\abs{x}$.
	\end{myproof}
\end{mytheorem}

Ein inkompressibles Wort erlaubt also keine Beschreibung, die kürzer ist als das Wort selbst. Es gibt somit keinen kürzeren Plan zu seiner Generierung als seine vollständige Beschreibung. Dies motiviert die folgende, von der gewählten Beschreibungssprache bis auf eine additive Konstante abhängige Definition.

\begin{mydefinition}[$c$-inkompressibles binäres Wort]\label{mydefinition:zufall}
\leavevmode
Es sei $c\geq0$ eine natürliche Zahl. Ein binäres Wort $w$ heisst $c$-inkompressibel, falls
\begin{align*}
	C(w)\geq\abs{w}-c.
\end{align*}
Ein $0$-inkompressibles Wort nennen wir kurz inkompressibel. Diese Begriffe beziehen sich stets auf die fest gewählte optimale universelle Beschreibungssprache.
\end{mydefinition}

\begin{myexercise}
	\leavevmode
	Warum bietet es sich bei einer natürlichen Zahl $n>0$ an, ihre Inkompressibilität relativ zur Länge von $bin(n)$ zu beurteilen?
	\begin{myanswer}
		Nach Definition gilt $C(n)=C(bin(n))$. Die Binärdarstellung besitzt die Länge $\abs{bin(n)}=\ceil{\log_2(n+1)}$. Da ihre erste Stelle für $n>0$ stets eine $1$ ist, enthält sie nur $\abs{bin(n)}-1$ frei wählbare Bits. Jede exakte Schwelle bleibt jedoch bis auf eine additive Konstante von der gewählten optimalen universellen Beschreibungssprache abhängig.
	\end{myanswer}
\end{myexercise}

\clearpage
\section{Die Kolmogorov-Komplexität ist nicht berechenbar}

\begin{mytheorem}[$C(w)$ ist nicht für jedes binäre Wort berechenbar]\label{mytheorem:Knotcomputable}
	Das Problem, für jedes binäre Wort $w$ die Kolmogorov-Komplexität $C(w)$ zu berechnen, ist algorithmisch unlösbar (nicht berechenbar).

	\vspace{\baselineskip}
	Mit anderen Worten: Es existiert kein Algorithmus, der ein beliebiges binäres Wort $w$ als Eingabe erhält und die natürliche Zahl $C(w)$ als Ausgabe liefert.
	\begin{myproof}
		Es sei $n>0$ eine natürliche Zahl. Wir bezeichnen mit $\alpha_n$ das erste Wort bezüglich der kanonischen Ordnung über dem binären Alphabet, das eine Kolmogorov-Komplexität von mindestens $n$ hat.
		Wegen \cref{mytheorem:nichtkomprimierbar} und der Verwendung der kanonischen Ordnung ist die Existenz und Eindeutigkeit von $\alpha_n$ gesichert.
		Beachte, dass wir keine Aussage zur Länge von $\alpha_n$ machen. Insbesondere wird nicht behauptet, dass $\alpha_n$ die Länge $n$ hat.

		\vspace{\baselineskip}
		Angenommen, es existiert ein Algorithmus $Q$, der zu jedem binären Wort $w$ die Zahl $C(w)$ berechnet. Dann können wir $Q$ verwenden, um eine kurze Beschreibung von $\alpha_n$ anzugeben. Für jede natürliche Zahl $n>0$ sucht und generiert das folgende \texttt{C'++}-Programm\footnote{Genauer: in \texttt{C'++}-Pseudocode.} das Wort $\alpha_n$:
		\begin{lstlisting}[language=C++]
			#include <iostream>
			#include <string>
			void seeking_alpha_n() {
				x = leeres Wort;
				Cx = Q(x);
				while (Cx < ***n+++) {
					x = Nachfolger von x in der kanonischen Ordnung über {0, 1}^*;
					Cx = Q(x);
				}
				std::cout << x;
			}
        \end{lstlisting}
		Alle Programme \cppinline{seeking_alpha_n} sind bis auf die Zahl $n$ identisch. Es sei $c$ die Länge ihres festen Maschinencodeteils. Damit gilt
		\begin{align*}
			C(\alpha_n)\leq c+\ceil{\log_2(n+1)}.
		\end{align*}

		\vspace{\baselineskip}
		Nach Definition von $\alpha_n$ gilt andererseits $C(\alpha_n)\geq n$. Zusammen erhalten wir für jedes $n>0$
		\begin{align*}
			n \leq C(\alpha_n) \leq c + \ceil{\log_2(n+1)}.
		\end{align*}
		Diese Ungleichung kann jedoch nur für endlich viele natürliche Zahlen $n>0$ gelten, weil $\log_2(n)$ langsamer wächst als $n$. Dies widerspricht der Annahme, dass ein Algorithmus $Q$ die Kolmogorov-Komplexität aller binären Wörter berechnet.
	\end{myproof}
\end{mytheorem}

\section{Elemente rekursiver Sprachen haben eine niedrige Kolmogorov-Komplexität}
TODO

\clearpage
\section{Verständnisaufgaben}

\begin{myexercise}
	Wie viele Stellen hat die kürzeste Darstellung der Zahl $n = 7^4 + 358$ in Basis $b = 6$?
	\begin{myanswer}
		Die gesuchte Länge ist
		\begin{align*}
			\ceil{\log_6\lr{7^4 + 358 + 1}} = \ceil{\ln\lr{7^4 + 359} / \ln\lr{6}} = 5.
		\end{align*}
	\end{myanswer}
\end{myexercise}

\begin{myexercise}
	Berechne
	\begin{align*}
		\sum_{k=0}^9 \lr{5 \cdot 4^k}.
	\end{align*}

	\begin{myanswer}
		\begin{align*}
			\sum_{k=0}^9 \lr{5 \cdot 4^k} &= 5\cdot \sum_{k=0}^9 \lr{4^k} = 5\cdot \frac{4^{10} - 1}{4 - 1} =\\
			&= 1747625
		\end{align*}
	\end{myanswer}
\end{myexercise}

\begin{myexercise}
		Es sei $q \neq 1$ eine reelle Zahl und es seien $m\leq n$ natürliche Zahlen. Dann gilt
	\begin{align*}
		\sum_{k = m}^n q^k = q^m + q^{m+1} + \ldots + q^n = \frac{q^{n+1}-q^m}{q-1}.
	\end{align*}
	\begin{myanswer}
		Siehe den Beweis von \cref{mytheorem:verallgemeinertegeometrischeSumme}.
	\end{myanswer}
\end{myexercise}

\begin{myexercise}
		Wie viele Elemente enthält die Menge
	\begin{align*}
		S := \setct{w}{$w$ ist ein binäres Wort mit $1\leq\abs{w}\leq n-1$}
	\end{align*}
	\begin{myanswer}
		Es gibt genau $2^1$ binäre Wörter der Länge $1$, $2^2$ binäre Wörter der Länge $2$ und allgemein $2^k$ binäre Wörter der Länge $k$. Insgesamt enthält die Menge $S$ also genau
		\begin{align*}
			\sum_{k = 1}^{n-1} 2^k \stackrel{\cref{mytheorem:verallgemeinertegeometrischeSumme}}{=} 2^n - 2
		\end{align*}
		Elemente.
	\end{myanswer}
\end{myexercise}
